\documentclass[sigconf,nonacm=false]{acmart}

% Remove ACM reference format for submission
\settopmatter{printacmref=false}
\renewcommand\footnotetextcopyrightpermission[1]{}

\usepackage{booktabs}
\usepackage{multirow}
\usepackage{graphicx}
\usepackage{amsmath}
\usepackage{xcolor}
\usepackage[normalem]{ulem} % for \sout (strikethrough)

\begin{document}

\title{Spiking Neural Networks for Environmental Sound Classification: From Seven Encodings to SpiNNaker Deployment}

\author{Ayush Kumar}
\affiliation{
  \institution{University of Manchester}
  \city{Manchester}
  \country{United Kingdom}
}
\email{ayush.kumar@student.manchester.ac.uk}

% Add supervisor as co-author (update name before submission)
%\author{Supervisor Name}
%\affiliation{
%  \institution{University of Manchester}
%}
%\email{supervisor@manchester.ac.uk}

%%% --- ABSTRACT (target: 150-200 words; current: ~175 words) ---
\begin{abstract}
Environmental sound classification (ESC) is critical for edge audio sensing, yet spiking neural networks (SNNs) remain untested on standard benchmarks.
We present the first convolutional SNN evaluation on ESC-50 (50 classes, 5-fold cross-validation), comparing seven spike encoding methods and deploying on SpiNNaker neuromorphic hardware.
Our SNN achieves 47.15\% $\pm$ 4.50\% with direct encoding---a 16.7 percentage point (pp) gap versus the matched-architecture ANN baseline (63.85\% $\pm$ 3.07\%, $p$=0.001).
This gap collapses to 0.95~pp ($p$=0.034) with frozen AudioSet-pretrained features, demonstrating the bottleneck is feature learning, not spiking computation.
SNNs exhibit 6.0$\times$ greater adversarial robustness under FGSM attack ($\varepsilon$=0.1, $p$=0.007, 5-fold validated).
Cross-encoding transfer analysis reveals a transfer ratio of 0.255 $\pm$ 0.006, showing encodings are highly task-specific with 75\% accuracy loss when swapped---a novel finding validated across all five folds.
We deploy on SpiNNaker via a validated FC2-only hybrid approach, achieving 33.1\% $\pm$ 6.9\% across five folds---the first neuromorphic hardware deployment for environmental sound classification.
\end{abstract}

\begin{CCSXML}
<ccs2012>
<concept>
<concept_id>10010147.10010178.10010179</concept_id>
<concept_desc>Computing methodologies~Neural networks</concept_desc>
<concept_significance>500</concept_significance>
</concept>
<concept>
<concept_id>10010583.10010588.10010591</concept_id>
<concept_desc>Hardware~Neural systems</concept_desc>
<concept_significance>500</concept_significance>
</concept>
<concept>
<concept_id>10010147.10010178.10010224.10010245.10010251</concept_id>
<concept_desc>Computing methodologies~Bio-inspired approaches</concept_desc>
<concept_significance>300</concept_significance>
</concept>
</ccs2012>
\end{CCSXML}

\ccsdesc[500]{Computing methodologies~Neural networks}
\ccsdesc[500]{Hardware~Neural systems}
\ccsdesc[300]{Computing methodologies~Bio-inspired approaches}

\keywords{spiking neural networks, environmental sound classification, neuromorphic computing, SpiNNaker, spike encoding}

\maketitle

%% ============================================================
\section{Introduction}
\label{sec:intro}
%% ============================================================

The proliferation of always-on audio sensing in edge devices demands both high accuracy and energy efficiency. Convolutional neural networks achieve near-human accuracy on environmental sound benchmarks~\cite{piczak2015}, but their multiply-accumulate (MAC) operations prohibit deployment on battery-constrained devices. Spiking neural networks (SNNs), which communicate via binary spike events, offer an alternative: on neuromorphic hardware such as SpiNNaker~\cite{furber2014} and Loihi~\cite{davies2018}, SNNs execute as accumulate-only operations (ACs) rather than MACs, with per-operation energy 5.1$\times$ lower~\cite{horowitz2014}.

Despite extensive SNN research on image classification~\cite{eshraghian2023}, \textbf{no prior work has evaluated convolutional SNNs on ESC-50}. The closest work~\cite{larroza2025} evaluates only ESC-10 (10 classes) with fully-connected layers. Dominguez-Morales et al.~\cite{dominguez2016} deploy SNNs on SpiNNaker for audio but evaluate only pure tones, not complex soundscapes.

This paper makes the following contributions:
\begin{enumerate}
    \item \textbf{First convolutional SNN on ESC-50} with systematic comparison of seven spike encoding methods, including a novel cross-encoding transfer analysis revealing high encoding specificity (transfer ratio = 0.255);
    \item \textbf{First SpiNNaker deployment for environmental sound}, using a validated FC2-only hybrid approach as proof of concept, with documented root-cause analysis of full-deployment failure;
    \item \textbf{Gap-collapse finding}: the SNN-ANN gap narrows from 16.7~pp to 0.95~pp with pre-trained features, showing the bottleneck is feature learning, not spiking computation;
    \item \textbf{First adversarial robustness analysis of SNNs on audio}, revealing 6.0$\times$ greater robustness under FGSM ($p$=0.007, 5-fold validated).
\end{enumerate}

%% ============================================================
\section{Background}
\label{sec:background}
%% ============================================================

\textbf{ESC-50}~\cite{piczak2015} contains 2{,}000 five-second recordings across 50 classes, with 5 predefined folds. Human performance is 81.3\%; ANN state of the art reaches 98--99\%~\cite{omnivec2}. We use log-mel spectrograms (64 bins, $n_\text{fft}$=1024, hop=512, sr=22050\,Hz), yielding 64$\times$216 input tensors.

\textbf{Spiking Neural Networks.} Leaky Integrate-and-Fire (LIF) neurons accumulate weighted input, spike when membrane potential exceeds a threshold, and reset. Gradients through the non-differentiable spike function use surrogate gradients~\cite{neftci2019,zenke2021}; we use fast sigmoid (slope=25). Training follows snnTorch~\cite{eshraghian2023} with cross-entropy on summed membrane potentials over $T$=25 timesteps.

\textbf{SpiNNaker}~\cite{furber2014} is a massively-parallel neuromorphic platform with integer weights and spike-driven communication. We use IF\_curr\_exp neurons with $\tau_\text{syn}$=5\,ms, $v_\text{thresh}$=1.0, $\tau_m$=20\,ms, calibrated via a 9-point scale sweep.

\textbf{Prior SNN Audio Work.} Larroza et al.~\cite{larroza2025} classify ESC-10 with FC-only SNNs. Dominguez-Morales et al.~\cite{dominguez2016} deploy SNNs on SpiNNaker for 8 pure tones. Neither addresses the full 50-class ESC-50 challenge or evaluates convolutional architectures.

%% ============================================================
\section{Methods}
\label{sec:methods}
%% ============================================================

\subsection{Architecture}

\begin{figure}[t]
\centering
\fbox{\parbox{0.93\columnwidth}{\centering\small\vspace{2pt}
\textbf{SpikingCNN} ($\sim$622K params, $T$=25)\\[4pt]
Input (1$\times$64$\times$216) \\
$\downarrow$ \\
Conv2d(1$\to$32, 3$\times$3) $\to$ BN $\to$ MaxPool(2) $\to$ LIF$_1$ \\
$\downarrow$ \\
Conv2d(32$\to$64, 3$\times$3) $\to$ BN $\to$ MaxPool(2) $\to$ LIF$_2$ \\
$\downarrow$ \\
AvgPool(4$\times$6) $\to$ Flatten \\
$\downarrow$ \\
FC(2304$\to$256) $\to$ LIF$_3$ $\to$ FC(256$\to$50) $\to$ LIF$_4$
\vspace{2pt}}}
\caption{SpikingCNN architecture. The ANN baseline replaces each LIF with ReLU. Both have $\sim$622K parameters.}
\label{fig:architecture}
\end{figure}

Both SNN and ANN share an identical architecture ($\sim$622K parameters, Figure~\ref{fig:architecture}). The ANN replaces each LIF with ReLU. Training uses Adam (lr=$10^{-3}$, wd=$10^{-4}$), ReduceLROnPlateau (factor=0.5, patience=5), early stopping (patience=10), 50 epochs, batch size 32, with standard 5-fold cross-validation.

\subsection{Spike Encoding Methods}

We evaluate seven encoding methods (Table~\ref{tab:encodings}), the most comprehensive comparison for audio SNNs to date. Encodings convert the 64$\times$216 spectrogram into spike trains across $T$=25 timesteps.

\begin{table}[t]
\caption{Seven spike encoding methods evaluated.}
\label{tab:encodings}
\small
\begin{tabular}{@{}lll@{}}
\toprule
Encoding & Description & Key Property \\
\midrule
Rate & Spike prob.\ $\propto$ intensity & Stochastic \\
Latency & High intensity $\to$ earlier spike & Temporal \\
Delta & Spikes on intensity changes & Sparse \\
Direct & Raw values repeated $T$ times & No conversion \\
Burst & Spikes concentrated at start & Dense, early \\
Phase & Spike time in oscillation cycle & Deterministic \\
Population & 10 neurons/class, MSE loss & Multi-neuron \\
\bottomrule
\end{tabular}
\end{table}

\subsection{SpiNNaker Deployment}

\begin{figure}[t]
\centering
\fbox{\parbox{0.93\columnwidth}{\centering\small\vspace{2pt}
\textbf{Hybrid SpiNNaker Pipeline}\\[4pt]
\textit{Software (snnTorch on CPU):}\\
Spectrogram $\to$ Conv+BN+Pool+LIF$_{1,2}$ $\to$ AvgPool $\to$ FC$_1$+LIF$_3$\\
$\downarrow$ \textit{binary hidden spikes (256-d, 21.7\% active/step)}\\[2pt]
\textit{SpiNNaker Hardware (FC$_2$-only):}\\
FC$_2$(256$\to$50) + LIF$_4$ $\to$ class prediction\\[4pt]
\sout{Full FC$_1$+FC$_2$ deployment} $\leftarrow$ \textit{failed: exc-inh cancellation}
\vspace{2pt}}}
\caption{Hybrid SpiNNaker deployment pipeline. Full deployment failed because AvgPool outputs are fractional, causing excitatory-inhibitory cancellation in FC$_1$ (see Section~\ref{sec:spinnaker_results}).}
\label{fig:pipeline}
\end{figure}

Full SNN deployment on SpiNNaker failed because AvgPool produces fractional (non-binary) outputs: FC$_1$ weights have near-zero mean ($\mu$=$-$0.0034), and 1{,}398 simultaneous fractional inputs produce net current $\approx$0, yielding zero hidden spikes. We adopt an FC2-only hybrid (Figure~\ref{fig:pipeline}): conv+FC$_1$+LIF$_3$ run in software, producing binary hidden spikes (256-d, 21.7\% active per step). Only FC$_2$ (256$\to$50)+LIF$_4$ runs on SpiNNaker. This constitutes a proof-of-concept deployment, not full neuromorphic inference.

\subsection{NeuroBench Energy Analysis}

We follow the NeuroBench framework~\cite{yik2025} to count effective ACs (SNN) and MACs (ANN). Energy estimation uses Horowitz ISSCC 2014 costs~\cite{horowitz2014}: AC=0.9\,pJ, MAC=4.6\,pJ.

\subsection{Adversarial Robustness Evaluation}

We apply FGSM~\cite{goodfellow2015} and PGD~\cite{madry2018} attacks across seven $\varepsilon$ values (0.00--0.20) using the torchattacks library. All results are 5-fold validated. Prior work on image classification~\cite{sharmin2020} shows SNNs exhibit inherent adversarial robustness from discrete encoding; we provide the first such analysis on audio spectrograms. We note that Wang et al.~\cite{wang2025} recommend SA-PGD for stronger SNN evaluation; our standard PGD results represent a lower bound on SNN robustness.

%% ============================================================
\section{Results}
\label{sec:results}
%% ============================================================

\subsection{Encoding Comparison}

\begin{table}[t]
\caption{ESC-50 accuracy (5-fold CV). All pairwise comparisons are statistically significant ($p<$0.002) except rate $\approx$ phase ($p$=0.93). Cohen's $d$ ranges from 3.75 to 8.13.}
\label{tab:encoding_results}
\small
\begin{tabular}{@{}lrrl@{}}
\toprule
Model / Encoding & Mean (\%) & $\pm$ Std & Significance \\
\midrule
\textbf{ANN baseline} & \textbf{63.85} & 3.07 & --- \\
\midrule
SNN Direct & \textbf{47.15} & 4.50 & $p$=0.001 vs ANN \\
SNN Phase & 24.15 & 1.66 & $p<$0.001 vs direct \\
SNN Rate & 24.00 & 1.90 & $p$=0.93 vs phase \\
SNN Population & 19.15 & 2.79 & $p<$0.002 vs rate \\
SNN Latency & 16.30 & 1.62 & $p<$0.002 vs pop. \\
SNN Delta & 7.25 & 0.94 & $p<$0.001 vs lat. \\
SNN Burst & 6.50 & 1.54 & $p<$0.002 vs delta \\
\bottomrule
\end{tabular}
\end{table}

\begin{figure}[t]
\centering
\fbox{\parbox{0.93\columnwidth}{\centering\small\vspace{2pt}
\textbf{Encoding Accuracy Comparison}\\[4pt]
\begin{tabular}{@{}lr@{}}
ANN baseline & \rule{60mm}{3mm} 63.85\% \\[1pt]
SNN Direct & \rule{44mm}{3mm} 47.15\% \\[1pt]
SNN Phase & \rule{18mm}{3mm} 24.15\% \\[1pt]
SNN Rate & \rule{18mm}{3mm} 24.00\% \\[1pt]
SNN Population & \rule{14mm}{3mm} 19.15\% \\[1pt]
SNN Latency & \rule{12mm}{3mm} 16.30\% \\[1pt]
SNN Delta & \rule{5mm}{3mm} 7.25\% \\[1pt]
SNN Burst & \rule{5mm}{3mm} 6.50\% \\
\end{tabular}
\vspace{2pt}}}
\caption{ESC-50 accuracy across seven SNN encodings and ANN baseline (5-fold CV). Direct encoding dominates all spike-converted methods by $>$20~pp. [Replace with generated bar chart for camera-ready.]}
\label{fig:encodings}
\end{figure}

Direct encoding outperforms all spike-converted alternatives by $>$20~pp (Table~\ref{tab:encoding_results}, Figure~\ref{fig:encodings}). The near-equality of rate (24.00\%) and phase (24.15\%, $p$=0.93) is notable: phase produces exactly one spike per neuron while rate produces $\sim$7, yet both achieve statistically identical accuracy. This confirms that temporal window coverage, not spike count, determines accuracy. The 16.7~pp SNN-ANN gap is statistically significant (paired $t$-test: $p$=0.001, Cohen's $d$=$-$2.93).

\textbf{Why do encodings fail?} Delta (7.25\%) fails because static spectrograms have no temporal variation. Burst (6.50\%) concentrates spikes in the first 5 of 25 timesteps, creating a temporal mismatch that causes severe overfitting (train 45--62\%, test 5--9\%).

\subsection{Encoding Transfer Analysis}

\begin{table}[t]
\caption{Cross-encoding transfer ratio (5-fold validated). Models trained with one encoding are evaluated using all others. Low ratios indicate high encoding specificity.}
\label{tab:transfer}
\small
\begin{tabular}{@{}lr@{}}
\toprule
Metric & Value \\
\midrule
Transfer ratio (mean $\pm$ std) & 0.255 $\pm$ 0.006 \\
Mean accuracy loss when swapped & 75\% \\
Fold-to-fold consistency ($n$=5) & $<$2.4\% variation \\
\bottomrule
\end{tabular}
\end{table}

We introduce a cross-encoding transfer analysis: models trained with one encoding are evaluated using all other encodings. The transfer ratio (accuracy with mismatched encoding divided by accuracy with matched encoding) averages 0.255~$\pm$~0.006 across all five folds (Table~\ref{tab:transfer}). This remarkably consistent result ($<$2.4\% coefficient of variation) demonstrates that SNN representations are highly encoding-specific---swapping the encoding at test time causes 75\% accuracy loss on average. To our knowledge, this is the first quantification of encoding specificity in SNNs on any benchmark.

\subsection{SpiNNaker Hardware Results}
\label{sec:spinnaker_results}

\begin{table}[t]
\caption{SpiNNaker 5-fold results (FC2-only hybrid proof of concept, 2{,}000 total inferences). Hardware gap is statistically significant ($p$=0.0016, paired $t$-test).}
\label{tab:spinnaker}
\small
\begin{tabular}{@{}lrrr@{}}
\toprule
Fold & SpiNNaker (\%) & snnTorch (\%) & Gap (pp) \\
\midrule
1 & 29.0 & 43.25 & 14.2 \\
2 & 32.0 & 44.0 & 12.0 \\
3 & 36.5 & 47.75 & 11.2 \\
4 & 43.0 & 51.25 & 8.2 \\
5 & 25.2 & 43.5 & 18.3 \\
\midrule
\textbf{Mean} & \textbf{33.1 $\pm$ 6.9} & \textbf{46.0} & \textbf{12.8 $\pm$ 4.1} \\
\bottomrule
\end{tabular}
\end{table}

Table~\ref{tab:spinnaker} presents the first neuromorphic hardware deployment for environmental sound classification. The 12.8~pp hardware gap ($p$=0.0016, paired $t$-test) arises from integer weight quantisation and IF neuron approximation. While accuracy is modest, the prior art~\cite{dominguez2016} evaluated only 8 pure tones; our 50-class task is fundamentally harder.

\textbf{Root-cause analysis of FC$_1$ failure.} Full SpiNNaker deployment requires binary inputs throughout. The AvgPool layer between LIF$_2$ and FC$_1$ produces fractional outputs in $[0,1]$. FC$_1$ weights have near-zero mean ($\mu$=$-$0.0034), so 1{,}398 simultaneous fractional inputs produce net excitatory-inhibitory cancellation ($I_\text{net} \approx 0$), yielding zero hidden spikes. Post-hoc weight re-centring (adding bias = $\mu \times n_\text{inputs}$) failed: accuracy dropped from 53.75\% to 8.50\% because the bias compensation assumes binary inputs but actual inputs are fractional. This reveals that standard conv$\to$pool$\to$FC architectures are not directly deployable on spike-only hardware without architectural redesign---a practically important co-design constraint.

\subsection{Transfer Learning: Gap Collapse}

\begin{table}[t]
\caption{PANNs CNN14 frozen embeddings + classifier head (5-fold CV). The SNN-ANN gap collapses from 16.7~pp to 0.95~pp.}
\label{tab:panns}
\small
\begin{tabular}{@{}lrrl@{}}
\toprule
Head & Mean (\%) & $\pm$ Std & vs SNN ($p$) \\
\midrule
PANNs + Linear & 93.80 & 1.69 & 0.011 \\
PANNs + ANN & 93.45 & 1.54 & 0.034 \\
PANNs + SNN & 92.50 & 1.30 & --- \\
\bottomrule
\end{tabular}
\end{table}

The SNN-ANN gap collapses from 16.7~pp (scratch) to 0.95~pp ($p$=0.034) with frozen AudioSet-pretrained CNN14 embeddings~\cite{kong2020} (Table~\ref{tab:panns}). Notably, the linear probe (93.80\%) outperforms both learned heads, indicating that CNN14 embeddings are already linearly separable for ESC-50. The SNN head ($p$=0.011 vs linear) adds computational overhead without accuracy benefit. Nevertheless, the key scientific insight holds: given equal-quality features, the SNN achieves accuracy within 0.95~pp of the ANN, demonstrating that the scratch-training gap is a \emph{feature-learning} problem, not a \emph{spiking computation} limitation. In a practical hybrid deployment, CNN14 extracts features in software while the SNN head classifies on neuromorphic hardware.

\subsection{Adversarial Robustness}

\begin{table}[t]
\caption{Accuracy (\%) under adversarial attack (5-fold mean $\pm$ std). Clean accuracy differs from Table~\ref{tab:encoding_results} due to BatchNorm running statistics across evaluation hardware; robustness \emph{ratios} are unaffected (see text).}
\label{tab:adversarial}
\small
\begin{tabular}{@{}lrrrr@{}}
\toprule
$\varepsilon$ & FGSM SNN & FGSM ANN & PGD SNN & PGD ANN \\
\midrule
0.00 & 54.25\tiny{$\pm$5.1} & 61.85\tiny{$\pm$4.6} & 54.25\tiny{$\pm$5.1} & 61.85\tiny{$\pm$4.6} \\
0.01 & 33.55\tiny{$\pm$3.7} & 21.05\tiny{$\pm$4.4} & 18.65 & 12.65 \\
0.02 & 29.30\tiny{$\pm$4.3} & 9.75\tiny{$\pm$2.2} & 13.65 & 2.05 \\
0.05 & 22.75\tiny{$\pm$4.1} & 5.05\tiny{$\pm$1.9} & 9.75 & 0.05 \\
\textbf{0.10} & \textbf{16.55}\tiny{$\pm$5.5} & \textbf{2.75}\tiny{$\pm$0.6} & \textbf{3.50} & \textbf{0.00} \\
0.20 & 13.25\tiny{$\pm$6.8} & 2.50\tiny{$\pm$0.7} & 1.15 & 0.00 \\
\bottomrule
\end{tabular}
\end{table}

Under FGSM at $\varepsilon$=0.1, the SNN retains 16.55\% $\pm$ 5.49\% accuracy versus 2.75\% $\pm$ 0.61\% for the ANN (\textbf{6.0$\times$ more robust}, $p$=0.007, 5-fold validated; Table~\ref{tab:adversarial}). Under PGD at $\varepsilon$=0.05, the SNN retains 9.75\% while the ANN drops to 0.05\%. Binary spike thresholding provides natural gradient masking: small perturbations that shift continuous activations have no effect until crossing the LIF threshold.

\textbf{Clean accuracy note.} The $\varepsilon$=0.00 values (SNN=54.25\%, ANN=61.85\%) differ from Table~\ref{tab:encoding_results} due to BatchNorm running statistics on different evaluation hardware; robustness \emph{ratios} are unaffected as both numerator and denominator use the same pipeline.

\subsection{Energy Analysis}

\begin{table}[t]
\caption{NeuroBench energy comparison (single fold). SNN activation sparsity = 74.2\% (spike rate = 25.8\%).}
\label{tab:energy}
\small
\begin{tabular}{@{}lrrr@{}}
\toprule
Model & Eff.\ Ops & Energy (nJ) & Op Type \\
\midrule
SNN & 1.08M ACs & 976 & AC $\times$ 0.9\,pJ \\
ANN & 101K MACs & 463 & MAC $\times$ 4.6\,pJ \\
\bottomrule
\end{tabular}
\end{table}

In software simulation, the ANN is 2.1$\times$ more energy-efficient (Table~\ref{tab:energy}) because the SNN's $T$=25 timesteps generate 10.7$\times$ more operations despite each being cheaper. On neuromorphic hardware, each AC costs 5.1$\times$ less than each MAC~\cite{horowitz2014}; however, the SNN's 25.8\% spike rate exceeds the $<$6.4\% threshold identified by Dampfhoffer et al.~\cite{dampfhoffer2023} for energy parity with quantised ANNs. Reducing $T$ and increasing sparsity are necessary for the SNN to achieve energy advantage---this is an important direction for future work.

%% ============================================================
\section{Discussion}
\label{sec:discussion}
%% ============================================================

\textbf{Why does direct encoding dominate?} Direct encoding feeds continuous values at every timestep; LIF neurons perform internal temporal integration via membrane dynamics. Spike-converted inputs lose information through discretisation: rate discards magnitude, latency collapses temporal richness, and delta amplifies noise on static inputs. Direct encoding is the recommended default for spectrogram-based SNN processing.

\textbf{The gap is feature-learning, not spiking.} The PANNs result (Table~\ref{tab:panns}) confirms: given equal-quality features, SNN $\approx$ ANN. The 16.7~pp scratch-training gap reflects difficulty learning convolutional filters from 1{,}600 samples with approximate surrogate gradients~\cite{neftci2019} and binary spike bottlenecks, consistent with Deng \& Gu~\cite{deng2020}. The linear probe (93.80\%) outperforming both heads confirms CNN14 features are already well-separated.

\textbf{Encoding specificity as a design constraint.} The transfer ratio of 0.255 means SNNs retain only 25.5\% accuracy with mismatched encodings. This high specificity---the SNN learns encoding-dependent representations---has practical implications for neuromorphic systems where the sensor encoding may be fixed by hardware.

\textbf{SpiNNaker: proof of concept with honest limitations.} Our FC2-only hybrid demonstrates feasibility but is not full neuromorphic inference. The FC$_1$ cancellation (Section~\ref{sec:spinnaker_results}) reveals that standard conv$\to$pool$\to$FC architectures with AvgPool are not directly deployable on spike-only hardware. Replacing AvgPool with MaxPool (which preserves binary outputs) could enable full deployment on SpiNNaker~2~\cite{hoppner2024}.

\textbf{Adversarial robustness.} Binary thresholding provides natural gradient masking, consistent with Sharmin et al.~\cite{sharmin2020} on images. We acknowledge that standard PGD may not be optimal for evaluating SNN robustness~\cite{wang2025}; SA-PGD evaluation is left for future work. Our FGSM results (single-step, no SNN-specific adaptation) provide a reliable lower bound.

\textbf{Threats to validity.} (1)~Results are for a single architecture ($\sim$622K params); wider/deeper networks may shift the encoding hierarchy. (2)~ESC-50 is small (2{,}000 clips); generalisation to UrbanSound8K or FSD50K is untested. (3)~NeuroBench energy is computed for one fold; actual SpiNNaker power consumption was not measured. (4)~The surrogate gradient (fast sigmoid) may not be optimal; preliminary single-fold ablation suggests spike-rate-escape achieves +1.25~pp.

%% ============================================================
\section{Conclusion}
\label{sec:conclusion}
%% ============================================================

We present the first convolutional SNN evaluation on ESC-50, establishing a 47.15\% baseline with direct encoding and deploying on SpiNNaker neuromorphic hardware as a proof of concept. The key finding is that the SNN-ANN gap is not fundamental---it collapses from 16.7~pp to 0.95~pp with pre-trained features ($p$=0.034). SNNs additionally exhibit 6.0$\times$ greater adversarial robustness under FGSM ($p$=0.007), and cross-encoding transfer analysis reveals high encoding specificity (transfer ratio = 0.255). Future work includes SpiNNaker~2~\cite{hoppner2024} deployment with binary-compatible architectures, sparsity optimisation to reach the $<$6.4\% spike-rate threshold for energy parity, and cross-dataset validation on UrbanSound8K.

%% ============================================================
\begin{acks}
This work was conducted as part of an undergraduate thesis at the University of Manchester. SpiNNaker hardware access was provided by the Advanced Processor Technologies group at the University of Manchester. Computational experiments were performed on the CSF3 high-performance computing facility at the University of Manchester.
\end{acks}

%% ============================================================
\bibliographystyle{ACM-Reference-Format}
\begin{thebibliography}{19}

\bibitem{piczak2015}
K.~J. Piczak.
\newblock ESC: Dataset for Environmental Sound Classification.
\newblock In \textit{Proc.\ 23rd ACM Int.\ Conf.\ Multimedia (MM'15)}, 1015--1018, 2015.

\bibitem{furber2014}
S.~B. Furber, F.~Galluppi, S.~Temple, and L.~A. Plana.
\newblock The SpiNNaker Project.
\newblock \textit{Proc.\ IEEE}, 102(5):652--665, 2014.

\bibitem{eshraghian2023}
J.~K. Eshraghian, M.~Ward, E.~O. Neftci, X.~Wang, G.~Liang, B.~Linares-Barranco, and W.~D. Lu.
\newblock Training Spiking Neural Networks Using Lessons From Deep Learning.
\newblock \textit{Proc.\ IEEE}, 111(9):1016--1054, 2023.

\bibitem{neftci2019}
E.~O. Neftci, H.~Mostafa, and F.~Zenke.
\newblock Surrogate Gradient Learning in Spiking Neural Networks.
\newblock \textit{IEEE Signal Processing Magazine}, 36(6):51--63, 2019.

\bibitem{zenke2021}
F.~Zenke and T.~P. Vogels.
\newblock The Remarkable Robustness of Surrogate Gradient Learning for Instilling Complex Function in Spiking Neural Networks.
\newblock \textit{Neural Computation}, 33(4):899--925, 2021.

\bibitem{dampfhoffer2023}
M.~Dampfhoffer, T.~Mesquida, A.~Valentian, and L.~Anghel.
\newblock Are SNNs Really More Energy-Efficient than ANNs? An In-Depth Hardware-Aware Study.
\newblock \textit{IEEE Trans.\ Emerging Topics in Computational Intelligence}, 7(3):731--741, 2023.

\bibitem{larroza2025}
A.~Larroza, L.~Reyes, J.~Maudes-Raedo, and G.~Rodriguez.
\newblock Evaluation of Spiking Neural Networks for Audio Classification.
\newblock \textit{arXiv:2503.11206}, 2025.

\bibitem{dominguez2016}
J.~P. Dominguez-Morales, \'{A}.~Jim\'{e}nez-Fern\'{a}ndez, A.~Rios-Navarro, E.~Cerezuela-Escudero, D.~Gutierrez-Galan, M.~J. Dom\'{i}nguez-Morales, and G.~Jim\'{e}nez-Moreno.
\newblock Multilayer Spiking Neural Network for Audio Samples Classification Using SpiNNaker.
\newblock In \textit{ICANN 2016}, LNCS 9886, pp.~45--53. Springer, 2016.

\bibitem{yik2025}
J.~Yik, S.~Ahmed, Z.~Ahmed, et~al.
\newblock NeuroBench: A Framework for Benchmarking Neuromorphic Computing Algorithms and Systems.
\newblock \textit{Nature Communications}, 16:1589, 2025.

\bibitem{kong2020}
Q.~Kong, Y.~Cao, T.~Iqbal, Y.~Wang, W.~Wang, and M.~D. Plumbley.
\newblock PANNs: Large-Scale Pretrained Audio Neural Networks for Audio Pattern Recognition.
\newblock \textit{IEEE/ACM Trans.\ Audio, Speech, and Language Processing}, 28:2880--2894, 2020.

\bibitem{sharmin2020}
S.~Sharmin, N.~Rathi, P.~Panda, and K.~Roy.
\newblock Inherent Adversarial Robustness of Deep Spiking Neural Networks: Effects of Discrete Input Encoding and Non-Linear Activations.
\newblock In \textit{ECCV 2020}, pp.~768--784, 2020.

\bibitem{goodfellow2015}
I.~J. Goodfellow, J.~Shlens, and C.~Szegedy.
\newblock Explaining and Harnessing Adversarial Examples.
\newblock In \textit{ICLR 2015}, 2015.

\bibitem{madry2018}
A.~Madry, A.~Makelov, L.~Schmidt, D.~Tsipras, and A.~Vladu.
\newblock Towards Deep Learning Models Resistant to Adversarial Attacks.
\newblock In \textit{ICLR 2018}, 2018.

\bibitem{omnivec2}
S.~Srivastava and A.~K. Sharma.
\newblock OmniVec2: A Novel Transformer Based Backbone for Perception Tasks.
\newblock In \textit{CVPR 2024}, 2024.

\bibitem{horowitz2014}
M.~Horowitz.
\newblock 1.1 Computing's Energy Problem (and What We Can Do About It).
\newblock In \textit{IEEE ISSCC Digest of Technical Papers}, pp.~10--14, 2014.

\bibitem{deng2020}
S.~Deng and S.~Gu.
\newblock Rethinking the Performance Comparison Between SNNs and ANNs.
\newblock \textit{Neural Networks}, 121:294--307, 2020.

\bibitem{wang2025}
J.~Wang, D.~Zhao, R.~Chen, Q.~Zhang, and Y.~Zeng.
\newblock Towards Reliable Evaluation of Adversarial Robustness for Spiking Neural Networks.
\newblock \textit{arXiv:2512.22522}, 2025.

\bibitem{hoppner2024}
S.~H\"{o}ppner, Y.~Yan, C.~Garbers, et~al.
\newblock SpiNNaker 2: A Large-Scale Neuromorphic System for Event-Based and Asynchronous Machine Learning.
\newblock \textit{arXiv:2401.04491}, 2024.

\bibitem{davies2018}
M.~Davies, N.~Srinivasa, T.-H. Lin, et~al.
\newblock Loihi: A Neuromorphic Manycore Processor with On-Chip Learning.
\newblock \textit{IEEE Micro}, 38(1):82--99, 2018.

\end{thebibliography}

\end{document}
